% ID: MAT-TRI-EQG-008
% Titolo: Equazioni goniometriche e periodicità (2)
% Disciplina: Matematica
% Area: trigonometria
% Argomento: equazioni_goniometriche
% Classe: Quarta secondaria di secondo grado
% Difficolta: 4
% Tipo: problema_numerico
% Tag: matematica, trigonometria, equazioni_goniometriche
% Fonte: verifica_goniometria_intermedia_anonimizzata
% Provenienza: materiale fornito dal docente; testo anonimizzato; soluzione redatta per Exergo

\begin{esercizio}
Risolvi le seguenti equazioni: 

\begin{enumerate}
\item $\displaystyle \cos 2x=\sin x;$

\item $\displaystyle \sin 2x+\sqrt{3}\cos x=0.$
\end{enumerate}
Salvo diversa indicazione, lavora in $\mathbb R$; angoli in radianti.
\end{esercizio}

\begin{soluzione}
a) $1-2\sin^2x=\sin x$ dà $\sin x=1/2$ o $-1$:
$x=\pi/6+2k\pi,5\pi/6+2k\pi,3\pi/2+2k\pi$.
b) $\cos x(2\sin x+\sqrt3)=0$ dà
$x=\pi/2+k\pi,4\pi/3+2k\pi,5\pi/3+2k\pi$. In tutte le famiglie di soluzioni $k\in\mathbb{Z}$.
\end{soluzione}
